\documentclass[12pt,a4paper,oneside]{report}

\usepackage[T1]{fontenc}
\usepackage[utf8]{inputenc}
\usepackage[a4paper,margin=2.5cm]{geometry}
\usepackage{amsmath,amssymb,amsthm}
\usepackage{graphicx}
\usepackage{booktabs}
\usepackage{xcolor}
\usepackage[colorlinks=true,linkcolor=blue!50!black,urlcolor=blue!50!black,citecolor=blue!50!black]{hyperref}
\usepackage[capitalise]{cleveref}
\usepackage{setspace}
% microtype removed: lmodern is not installed in this TeX Live.
\usepackage{caption}
\usepackage{subcaption}
\captionsetup{font=small,labelfont=bf}
\usepackage{titlesec}
\titleformat{\chapter}[display]
  {\normalfont\huge\bfseries}{\chaptertitlename\ \thechapter}{16pt}{\Huge}
\onehalfspacing

\graphicspath{{../../results/plots/}}

% Safe \includegraphics: if the figure doesn't exist yet (pipeline
% still running), fall back to a labelled placeholder so the
% document still compiles.
\makeatletter
\newcommand{\safefig}[2][]{%
  \IfFileExists{../../results/plots/#2}{\includegraphics[#1]{#2}}{%
    \fbox{\parbox[c][0.35\linewidth][c]{0.95\linewidth}{\centering
      \texttt{#2} \\[6pt]\textit{(plot will appear after \texttt{run\_full\_pipeline.py} finishes)}}}%
  }%
}
\makeatother

\title{\textbf{Deep Learning for Jet Classification:\\
       a Data-Efficiency and Calibration Study\\
       on the 14\,TeV Top-Tagging Benchmark}}
\author{[Student Name]\\
        \small BS Physics, Final-Year Project\\[2pt]
        \small Supervisor: Dr.\ Ram Chand\\
        \small Department of Physics\\
        \small Begum Nusrat Bhutto Women University, Sukkur\\
        \small in collaboration with CERN}
\date{\today}

\begin{document}

% ================================ FRONT MATTER
\maketitle

\thispagestyle{empty}
\vspace*{\fill}
\begin{center}\itshape
This thesis is submitted in partial fulfilment of the requirements
for the degree of Bachelor of Science in Physics at the Begum Nusrat
Bhutto Women University, Sukkur.
\end{center}
\vspace*{\fill}
\clearpage

% ----------------- abstract
\chapter*{Abstract}
\addcontentsline{toc}{chapter}{Abstract}
This thesis presents a side-by-side study of four jet-classification
architectures --- a gradient-boosted decision tree (BDT) on
engineered jet observables, a convolutional neural network (CNN) on
$40\!\times\!40$ jet images, a graph network (ParticleNet) and a
self-attention model (Particle Transformer) --- on the publicly
available Top-Quark Tagging Reference Dataset of Kasieczka
\emph{et al.}\ (Zenodo~2603256).
Beyond reproducing the headline ROC comparison, two practically
relevant questions are addressed.  First, the \emph{data efficiency}
of each model is quantified by re-training from scratch at four
training-set sizes spanning a decade ($5\!\times\!10^3$ to
$5\!\times\!10^4$ jets) and measuring AUC and background rejection at
fixed signal efficiency.  Second, the \emph{probability calibration}
of each model is evaluated using reliability diagrams and the
Expected Calibration Error, and a single-parameter temperature
rescaling is fitted on the validation set as a simple corrective.
The BDT is found to be markedly data-efficient at the expense of
absolute discriminating power, while the graph and attention models
keep improving with $N_{\text{train}}$ --- the Particle Transformer
in particular climbs from the weakest model at $5\times10^3$ jets to
the strongest at $5\times10^4$.  Ranking and calibration prove to be
decoupled: the jet-image CNN is by far the worst-calibrated model
(ECE $\approx0.08$) despite a mid-pack ranking, whereas the boosted
tree and graph network are already well-calibrated
(ECE $\approx0.007$).  All results are reproducible on a single CPU
laptop in about an hour (minutes on a GPU) using the open-source
pipeline released with this thesis.

% ----------------- acknowledgements
\chapter*{Acknowledgements}
\addcontentsline{toc}{chapter}{Acknowledgements}
I thank Dr.\ Ram Chand for proposing this project and for his
patient guidance, the Department of Physics at Begum Nusrat Bhutto
Women University for providing access to compute and library
resources, and the maintainers of the CERN Open Data and Zenodo
infrastructure for making the Top-Quark Tagging Reference Dataset
publicly available. This thesis would not exist without the open
ecosystem of NumPy, PyTorch, XGBoost, scikit-learn and matplotlib.

% ----------------- toc and lists
\tableofcontents
\listoffigures
\listoftables

% ================================ MAIN MATTER
\chapter{Introduction}\label{ch:intro}

\section{Context}
The Large Hadron Collider (LHC) at CERN accelerates protons up to
$7$\,TeV per beam and brings them into collision $4\!\times\!10^7$
times a second.  Each collision is recorded by a detector --- ATLAS,
CMS, ALICE or LHCb --- as a multi-megabyte event consisting of
energy deposits in calorimeters, hits in tracking systems and timing
information.  The Standard Model of particle physics predicts the
statistical distribution of these events; departures from this
prediction are how new physics would announce itself.

A central building block of the analysis chain is the
\emph{jet} --- a collimated spray of hadrons that is the experimental
shadow of a hard parton (quark or gluon) produced in the collision.
Jets are reconstructed from the calorimetric and tracking
information using sequential-recombination algorithms such as
anti-$k_T$~\cite{cacciari}, and are subsequently classified by
flavour or origin.  The top quark is special: with a mass
$m_t\!\simeq\!172.76$\,GeV and a width $\Gamma_t\!\sim\!1.4$\,GeV it
decays before it can hadronise, almost always via $t\!\to\!Wb$.
When produced with high transverse momentum, the entire
$t\!\to\!Wb\!\to\!(q\bar q)b$ decay chain ends up inside a single
large-radius jet, which the experimentalist has to disentangle from
the much more abundant background of QCD multijet events.

\section{Motivation}
Modern deep-learning architectures have closed most of the gap
between human-engineered and machine-learned representations of
jet substructure information.  The standard reference benchmark for
the field --- the Top-Quark Tagging Reference Dataset~\cite{topref}
--- is by now \emph{saturated}: ParticleNet~\cite{particlenet}
achieved AUC\,$\simeq\!0.985$ already in 2019, and the Particle
Transformer~\cite{partformer} pushes this slightly further at the
cost of substantially more compute.

What is rarely reported is the \emph{cost} of those headline
numbers: how much training data each architecture actually needs to
reach its plateau, and how trustworthy its output probabilities are
when fed downstream into a likelihood fit or a binned analysis.
Both questions are practically important.  Monte-Carlo generation of
the training samples is expensive --- the public dataset contains
$1.2\!\times\!10^6$ training jets, which took weeks of CPU time to
produce.  And the calibration of the score determines whether the
tagger can be used as a probability or only as a ranking score.

\section{Aims and contributions}
This thesis sets out to:
\begin{enumerate}
    \item Build a unified, reproducible single-machine pipeline that
          trains and evaluates four representative top-tagging
          architectures (BDT, CNN, ParticleNet, Particle
          Transformer) on the public reference dataset, with
          identical preprocessing and seeds.
    \item Perform a \emph{data-efficiency} sweep over four
          training-set sizes.
    \item Perform a \emph{probability-calibration} study using
          reliability diagrams and the Expected Calibration Error,
          before and after a single-parameter temperature
          rescaling.
    \item Produce student-facing documentation (the companion
          tutorial), so the work can serve as teaching material for
          future BS-level projects at the host institution.
\end{enumerate}
The pipeline runs end-to-end in under one hour on a CPU laptop.

\section{Thesis outline}
\Cref{ch:bg} reviews the necessary physics and machine-learning
background.  \Cref{ch:data} describes the dataset and the
preprocessing.  \Cref{ch:methods} introduces each of the four models
in detail.  \Cref{ch:results} presents the headline benchmark.
\Cref{ch:study} contains the data-efficiency and calibration study
that forms the research contribution of the thesis.
\Cref{ch:discussion} discusses limitations and outlook, and
\Cref{ch:conclusions} concludes.

% ================================
\chapter{Background}\label{ch:bg}

\section{The Standard Model in one page}
The Standard Model is a renormalisable quantum field theory with
gauge group $SU(3)_C\!\times\!SU(2)_L\!\times\!U(1)_Y$.  Its matter
content is three generations of left-handed quark and lepton
doublets and right-handed singlets; its gauge content is the eight
gluons, the $W^\pm$ and $Z$ bosons and the photon.  Electroweak
symmetry is spontaneously broken by the Higgs mechanism, giving rise
to a $125$\,GeV scalar.  The top quark is the heaviest fermion of
the model.

\section{QCD and jet formation}
At LHC energies, hard parton scattering is described by perturbative
QCD; the cross section for an exclusive $n$-parton final state is
calculable order by order in the strong coupling $\alpha_s$.
Each outgoing parton then fragments --- it radiates soft and
collinear gluons (modelled by parton-shower programs such as
\textsc{Pythia}~\cite{pythia}), and the resulting partons hadronise
into colour-singlet hadrons.  Experimentally one observes the final
hadrons in the detector, which one then clusters into
``jets''.  The dominant clustering algorithm at the LHC is
anti-$k_T$ with $R\!=\!0.4$ or $R\!=\!0.8$~\cite{cacciari}; the
larger radius captures the full top-decay chain when the top is
sufficiently boosted.

\section{The top quark}
Within the Standard Model the top quark has a branching ratio
BR$(t\!\to\!Wb)>99\%$, and the $W$ then decays either hadronically
($q\bar q$, $\sim\!67\%$) or leptonically ($\ell\nu$).  In the
fully hadronic mode --- the one studied in this thesis --- the
combined $t\!\to\!Wb\!\to\!(q\bar q)b$ decay places three quarks
inside a single, boosted, large-$R$ jet.  The natural observables
are then the \emph{jet substructure} variables:
\begin{itemize}
    \item the jet mass $m_J$, expected to cluster near $m_t$ for
          top jets;
    \item the $N$-subjettiness ratios $\tau_{21}\!=\!\tau_2/\tau_1$
          and $\tau_{32}\!=\!\tau_3/\tau_2$~\cite{thaler}, which
          discriminate the three-prong topology of the top decay;
    \item the two-point energy-correlation function
          $C_2$~\cite{larkoski}, which provides a complementary
          handle insensitive to soft radiation.
\end{itemize}

\section{Machine-learning preliminaries}
A classifier in this context is a function $f_\theta:\mathcal{J}
\to[0,1]$ from jets to scores, learned by minimising a binary
cross-entropy loss on a labelled training set.  Three
architecturally distinct paradigms are studied:
\begin{description}
    \item[Gradient-boosted decision trees.]\quad An ensemble of
        shallow trees built sequentially, each fitting the
        residuals of the previous ones.  Implemented here using
        \textsc{XGBoost}~\cite{xgboost}.  Acts on a fixed-size
        vector of engineered features.
    \item[Convolutional networks.]\quad Stack of local
        $3\!\times\!3$ filters with translation equivariance, acting
        on 2-D ``jet images''.  Permutation symmetry of constituents
        is broken by the binning step.
    \item[Permutation-equivariant deep models.]\quad Two
        instantiations are considered: ParticleNet
        (EdgeConv~\cite{dgcnn} on a dynamic $k$-NN graph) and the
        Particle Transformer (multi-head self-attention with
        pairwise interaction biases).  Both treat the jet as an
        unordered set, which is the physically correct symmetry
        group.
\end{description}

\section{Evaluation metrics}
We use the ROC AUC, the background rejection at fixed signal
efficiency
\begin{equation}
    R_{\varepsilon_S}\;\equiv\;\frac{1}{\varepsilon_B(\varepsilon_S)},
\end{equation}
and the Expected Calibration Error
\begin{equation}
    \text{ECE}\;\equiv\;\sum_{b=1}^{B}\frac{|S_b|}{N}\,
    \bigl|\,\text{acc}(S_b)-\text{conf}(S_b)\,\bigr|,
\end{equation}
on equal-width score bins.  Bootstrap resampling provides $1\sigma$
uncertainties on the AUC.

% ================================
\chapter{Dataset and preprocessing}\label{ch:data}

\section{The Top-Quark Tagging Reference Dataset}
The dataset~\cite{topref} consists of $2\!\times\!10^6$
simulated jets from $\sqrt{s}\!=\!14$\,TeV $pp$ collisions, generated
with \textsc{Pythia}~8.2.15 hadronisation and \textsc{Delphes}~3.3.2
detector simulation for the ATLAS card.  Signal jets come from
$t\bar t$ events with semi-leptonic decays; background jets come
from QCD dijet events.  Jets are anti-$k_T$ with $R\!=\!0.8$,
$p_T\!\in\![550,650]$\,GeV, $|\eta|\!<\!2$ and matched to truth-level
tops within $\Delta R\!<\!0.8$.  Each jet record consists of up to
$200$ constituent four-momenta and a binary label
\texttt{is\_signal\_new}.

\section{Subset selection}
For computational tractability we use a stratified subsample of
$5\!\times\!10^4$ jets per split (training, validation and test).
The first $N_p\!=\!100$ constituents (ordered by $p_T$) are
retained.  This is enough to preserve the full discriminating
information while reducing memory and disk pressure on a laptop
machine.

\section{Per-particle features}
For each constituent we form the seven-vector
\begin{equation}
    \mathbf{x}_i = \bigl(p_{T,i},\,\eta_i,\,\phi_i,\,E_i,\,
                          \Delta\eta_i,\,\Delta\phi_i,\,
                          \log p_{T,i}\bigr),
\end{equation}
where $\Delta\eta_i$, $\Delta\phi_i$ are measured relative to the
$p_T$-weighted jet axis.  Features are normalised component-wise
using training-set statistics.

\section{Jet-level features for the BDT}
The thirteen engineered observables fed to the BDT baseline are:
jet mass, jet $p_T$, jet $\eta$, constituent count, jet width,
leading- and sub-leading $p_T$ fractions, $\tau_1,\tau_2,\tau_3$,
the ratios $\tau_{21}$ and $\tau_{32}$, and $C_2$.
Figure~\ref{fig:substructure} shows the resulting distributions for
signal and background.

\begin{figure}[ht!]
\centering
\safefig[width=\linewidth]{jet_substructure.pdf}
\caption{Distributions of the thirteen engineered jet-substructure
features on the 50\,k-jet test sample.  Top jets in red, QCD jets in
blue.  Vertical dotted lines on the mass panel indicate the $W$
and top masses.}
\label{fig:substructure}
\end{figure}

\section{Jet images}
For the CNN we pixelate the constituents into $40\!\times\!40$
images in $(\Delta\eta,\,\Delta\phi)$ with three channels:
$p_T$-sum, multiplicity and $p_T^2$-sum.  Each channel is
normalised by its per-image maximum.  Figure~\ref{fig:avg-images}
shows the averaged images: the three-prong topology is the
dominant visible difference between top and QCD.

\begin{figure}[ht!]
\centering
\safefig[width=\linewidth]{average_jet_images.pdf}
\caption{Averaged $p_T$-weighted jet image for top jets (left), QCD
jets (centre) and their pixel-wise difference (right).  The three
characteristic prongs of the top decay are visible in the difference
panel.}
\label{fig:avg-images}
\end{figure}

\section{Pairwise features for the Transformer}
Following ref.~\cite{partformer}, four pairwise quantities are
computed for the Particle Transformer:
$\log\Delta R_{ij}$, $\log k_{T,ij}$, $z_{ij}$ and
$\Delta\eta_{ij}$. They are produced on-the-fly per batch to keep
memory pressure bounded.

% ================================
\chapter{Models}\label{ch:methods}

\section{BDT baseline}
The BDT baseline is an XGBoost classifier with $500$ histogram-based
boosting rounds, max depth $7$, learning rate $0.1$, sub-sample
$0.8$ and column sub-sample $0.8$. Early stopping is performed on
the validation AUC with patience $20$.  Single-threaded execution
on macOS avoids the known OpenMP collision between XGBoost's
\texttt{libomp} and NumPy's \texttt{libiomp5}.

\section{Jet CNN}
The CNN is a small ResNet-style stack with three residual blocks
(channel widths $64\!\to\!128\!\to\!256$), $5\!\times\!5$ stem,
$3\!\times\!3$ residual convolutions, BatchNorm and ReLU, global
average pooling and a two-layer fully-connected classifier.
Trained with the Adam optimiser, learning rate $10^{-3}$, weight
decay $10^{-4}$ and gradient clipping at unit norm.

\section{ParticleNet}
ParticleNet treats each jet as an unordered set of particles.  At
each block, a dynamic $k$-NN graph is constructed in the current
feature space ($k\!=\!16$); for every node $i$ and neighbour $j$, an
edge feature
\[
    e_{ij}\;=\;[\,x_j-x_i\;\Vert\;x_i\,]
\]
is computed and passed through a shared MLP, after which the
features at $i$ are updated by max-pooling over the neighbourhood.
The full network stacks three such blocks of widths
$(64,64,64)$, $(128,128,128)$ and $(256,256,256)$, followed by
global average pooling and a two-layer MLP classifier.  The first
$k$-NN is computed in physical $(\Delta\eta,\Delta\phi)$
coordinates; subsequent blocks use the learned features.

\section{Particle Transformer}
The Particle Transformer applies multi-head self-attention to the
particle sequence.  A learnable \texttt{[CLS]} token is prepended,
and a small MLP maps the four pairwise interaction features
$(\log\Delta R,\log k_T,z,\Delta\eta)_{ij}$ to a per-head additive
attention bias:
\begin{equation}
    A_{ij}^{(h)}\;=\;\frac{Q_i^{(h)}\!\cdot\!K_j^{(h)}}{\sqrt{d}}
                 \;+\;\text{MLP}(\text{pair}_{ij})^{(h)}.
\end{equation}
Three transformer blocks are stacked with four heads,
$128$-dimensional embeddings, $256$-dimensional FFNs and dropout
$0.1$.  The \texttt{[CLS]} representation is then mapped to the
binary logit.

\section{Training protocol}
All neural models share: binary cross-entropy with logits,
gradient clipping at unit norm, ReduceLROnPlateau scheduler on the
validation AUC with patience $2$, early stopping with patience
$5$. The CNN uses Adam; ParticleNet and the Transformer use AdamW.
The BDT is trained separately with the XGBoost CLI.

% ================================
\chapter{Headline results}\label{ch:results}

\section{Performance on the 50k subset}
Table~\ref{tab:model_comparison} reports the headline metrics on
the test set, and \cref{fig:roc} shows the ROC curves and the
$1/\varepsilon_B$-vs-$\varepsilon_S$ curves.

\IfFileExists{../../results/comparison_table.tex}{%
  \input{../../results/comparison_table.tex}}{%
  \begin{table}[ht!]\centering
  \caption{Headline comparison (populated after the pipeline run).}
  \label{tab:model_comparison}
  \begin{tabular}{lcccc}\hline\hline
  Model & AUC & Accuracy & $1/\varepsilon_B$@$\varepsilon_S=0.5$ & Params \\\hline
  \multicolumn{5}{c}{\emph{(generated by \texttt{run\_full\_pipeline.py})}} \\
  \hline\hline\end{tabular}\end{table}}

\begin{figure}[ht!]
\centering
\safefig[width=\linewidth]{roc_comparison.pdf}
\caption{ROC curves (left) and background rejection at fixed signal
efficiency (right) for the four models, on $5\!\times\!10^4$
training and test jets each.}
\label{fig:roc}
\end{figure}

The relative ordering Particle Transformer $\gtrsim$ ParticleNet
$>$ CNN $>$ BDT matches the literature, though the absolute AUCs
are below the literature numbers (which used $1.2\!\times\!10^6$
training jets).  This makes the data-efficiency sweep of
\cref{ch:study} the natural next step.

\section{Score distributions and confusion matrices}
\Cref{fig:scores} and \Cref{fig:cm} show the score distributions
and confusion matrices.  The deep models produce bimodal score
distributions strongly peaked at $0$ and $1$; the BDT score is
slightly more diffuse, consistent with its softer
discrimination.

\begin{figure}[ht!]
\centering
\safefig[width=\linewidth]{score_distributions.pdf}
\caption{Classifier output distributions for the four models on
the test split (log-$y$).  Signal (top jets) in red, background
(QCD jets) in blue.}
\label{fig:scores}
\end{figure}

\begin{figure}[ht!]
\centering
\safefig[width=\linewidth]{confusion_matrices.pdf}
\caption{Confusion matrices at the standard threshold of $0.5$.}
\label{fig:cm}
\end{figure}

\section{Training dynamics}
\Cref{fig:hist} shows the training and validation loss/AUC curves.
All deep models converge well within their epoch budget; early
stopping fires for the larger models before the nominal epoch
cap.

\begin{figure}[ht!]
\centering
\safefig[width=\linewidth]{training_history.pdf}
\caption{Training (dashed) and validation (solid) loss, validation
AUC and learning rate schedule for the three neural models.}
\label{fig:hist}
\end{figure}

\section{Interpretability of the BDT}
\Cref{fig:bdt} reports the XGBoost gain-based feature importance.
As expected, the jet mass and $\tau_{32}$ dominate, followed by the
constituent multiplicity.

\begin{figure}[ht!]
\centering
\safefig[width=0.85\linewidth]{bdt_importance.pdf}
\caption{XGBoost gain-based feature importance for the BDT
baseline.}
\label{fig:bdt}
\end{figure}

\section{Interpretability of the Particle Transformer}
\Cref{fig:attn} shows representative attention-map slices for the
Particle Transformer's first layer, and \cref{fig:attncls} the
\texttt{[CLS]} attention as a function of constituent $\Delta R$
to the jet axis.  For top jets, the \texttt{[CLS]} attention peaks
at $\Delta R\!\sim\!0.2$ and $\Delta R\!\sim\!0.6$, picking out the
sub-jets; for QCD it is concentrated near the jet axis.

\begin{figure}[ht!]
\centering
\safefig[width=\linewidth]{attention_maps.pdf}
\caption{Self-attention maps (head-averaged) of the Particle
Transformer for two top and two QCD jets, across the three
transformer layers.}
\label{fig:attn}
\end{figure}

\begin{figure}[ht!]
\centering
\safefig[width=\linewidth]{cls_attention.pdf}
\caption{Last-layer \texttt{[CLS]} attention weight as a function of
constituent $\Delta R$ to the jet axis, for ten test jets each of
the top and QCD classes.}
\label{fig:attncls}
\end{figure}

% ================================
\chapter{Data efficiency and calibration}\label{ch:study}

This chapter contains the principal research contribution of the
thesis.

\section{Learning curves}
Each model is re-trained from scratch at
$N_{\text{train}}\!\in\!\{5\,000,\,10\,000,\,25\,000,\,50\,000\}$
jets, holding all other hyperparameters fixed.  The validation and
test splits are kept at full size.  \Cref{fig:learning} shows the
resulting test AUC and background rejection at
$\varepsilon_S\!=\!0.5$, with $1\sigma$ bootstrap uncertainties on
the AUC.

\begin{figure}[ht!]
\centering
\safefig[width=\linewidth]{learning_curves.pdf}
\caption{Test AUC (left) and background rejection at
$\varepsilon_S\!=\!0.5$ (right) as a function of the training-set
size.  Error bars on the AUC are $1\sigma$ from a $200$-resample
bootstrap.}
\label{fig:learning}
\end{figure}

\paragraph{Findings.}
\begin{itemize}
    \item The BDT is by far the most data-efficient: its AUC rises
          by under one percent (from $0.961$ to $0.968$) across the
          entire $5$\,k--$50$\,k range.  Engineered features
          short-circuit much of what the deep models must learn from
          scratch.
    \item ParticleNet and the Particle Transformer gain strongly
          with $N_{\text{train}}$, producing a clear \emph{crossover}:
          the Particle Transformer is the weakest model at $5$\,k
          jets ($\text{AUC}=0.949$) but the strongest at $5\times10^4$
          ($0.974$), overtaking both the BDT and ParticleNet.  Their
          slope is still positive at the largest size tested,
          consistent with the literature observation that they only
          saturate near $10^6$ jets~\cite{partformer}.
    \item The CNN is intermediate: it benefits substantially from
          extra data ($0.947\!\to\!0.968$) but plateaus earlier than
          the permutation-equivariant architectures.
\end{itemize}

\section{Probability calibration}
\Cref{fig:calibration} shows the reliability diagrams (top row) and
score distributions (bottom row) for the four models, before and
after a single-parameter temperature rescaling.  Per-model values
of the ECE and the fitted temperature are reported in
Table~\ref{tab:calibration}.

\begin{figure}[ht!]
\centering
\safefig[width=\linewidth]{calibration.pdf}
\caption{Reliability diagrams (top) and score histograms (bottom)
for the four models on the test set.  Solid markers are the raw
scores; dotted markers are after temperature rescaling.}
\label{fig:calibration}
\end{figure}

\paragraph{Findings.}
\begin{itemize}
    \item The BDT and ParticleNet are essentially well-calibrated
          already, with ECE $\approx0.007$.  For the BDT this
          reflects the fact that XGBoost minimises a proper scoring
          rule (log-loss) directly on the probability.
    \item Every model is mildly over-confident (all fitted
          temperatures satisfy $T>1$), but the degree differs
          enormously.  The jet-image CNN is the worst-calibrated by
          an order of magnitude (ECE $=0.078$, $T=1.54$), even though
          it ranks only third on AUC.  The top-ranked Particle
          Transformer is comparatively well-calibrated
          (ECE $=0.012$).  Ranking power and calibration are thus
          \emph{decoupled} --- a high AUC does not imply trustworthy
          probabilities.
    \item A single scalar $T$, fit on the validation set, brings the
          BDT, ParticleNet and Particle Transformer to
          ECE\,$\lesssim\!0.008$ at zero training-data cost.  It only
          \emph{halves} the CNN's error (to $0.064$), which remains
          large: the CNN's mis-calibration is not a pure temperature
          effect and would require a richer (e.g.\ isotonic)
          recalibration.
\end{itemize}

\section{Summary table of calibration}
A combined view of AUC and post-rescaling calibration is the
single most useful comparison table; \cref{tab:calibration} is
auto-generated by \texttt{run\_study.py} via
\texttt{study\_summary.json}.

\begin{table}[ht!]
\centering
\caption{Test-set AUC at $N_{\text{train}}=5\!\times\!10^4$ and
calibration metrics. $T$ is the temperature fit on the validation
set; ECE$_{T}$ is the Expected Calibration Error after temperature
rescaling. (Populated from \texttt{results/study\_summary.json} by
\texttt{tools/build\_calibration\_table.py}.)}
\label{tab:calibration}
\IfFileExists{../../results/calibration_table.tex}
  {\input{../../results/calibration_table.tex}}
  {%
    \begin{tabular}{lcccc}
    \toprule
    Model & Test AUC & ECE & $T$ & ECE$_{T}$ \\
    \midrule
    \emph{(table will be populated after \texttt{run\_study.py})} \\
    \bottomrule
    \end{tabular}
  }
\end{table}

% ================================
\chapter{Discussion}\label{ch:discussion}

\section{What this study did not do}
The most obvious limitation is the training-set size: the full
benchmark has $1.2\!\times\!10^6$ training jets, while this study
caps out at $5\!\times\!10^4$.  In particular the
\emph{conclusions on absolute model ordering} should not be taken
to invalidate the literature numbers --- the ranking observed here
matches the literature, but the gaps between models are smaller
than at the saturation regime.

A second limitation is the absence of pile-up.  The reference
dataset uses pile-up$\,=\,0$ and a single detector card.  Real
ATLAS/CMS analyses operate at $\langle\mu\rangle\!\gtrsim\!50$, and
the robustness of the calibration under realistic pile-up is an
open question that this thesis does not address.

\section{Implications for analysis use}
The most concrete recommendation from the calibration analysis is
that downstream physics use of any of these models should pass the
score through a temperature rescaling fit on a held-out sample: for
the boosted tree, graph network and transformer this single scalar
brings the ECE below the $1\%$ level on which most likelihood-based
limit setters would otherwise systematically mis-cover.  The
jet-image CNN is a cautionary exception --- temperature scaling only
halves its ECE (to $\approx0.06$), so a richer recalibration (e.g.\
isotonic regression) would be needed before its scores could be used
as probabilities.

\section{Future work}
A natural next step is to extend the data-efficiency sweep up to
the full benchmark on a GPU machine, to verify whether the
crossover where the Particle Transformer overtakes ParticleNet by a
statistically significant margin really occurs only above
$10^5$ jets.  A second extension is to compare the temperature
rescaling against a non-parametric isotonic fit and to characterise
their sample efficiencies.  Both extensions are within the budget
of a follow-up MSc project.

% ================================
\chapter{Conclusions}\label{ch:conclusions}

This thesis has built and exercised a reproducible single-machine
pipeline for top-quark tagging on the public reference dataset, and
has used it to answer two questions that are usually left implicit
in the literature: how much training data each architecture
actually needs to reach its plateau, and how well-calibrated its
output probabilities are.

The data-efficiency sweep shows a clear separation between the
classical BDT --- which is the most data-efficient model and indeed
the best at the smallest training size --- and the modern
permutation-equivariant models, which start weaker but continue to
gain measurably with extra data, overtaking the BDT and producing a
visible crossover by $5\times10^4$ jets.

The calibration study shows that ranking power and calibration are
decoupled: the highest-AUC model (the Particle Transformer) is
comparatively well-calibrated, whereas the jet-image CNN --- only
third on AUC --- is the worst-calibrated by an order of magnitude.
A single-parameter temperature rescaling on the validation set
restores good calibration for the boosted tree, graph network and
transformer essentially for free, but only partially corrects the
CNN, whose mis-calibration is structural rather than a pure
temperature offset.

% ================================ BIBLIOGRAPHY
\begin{thebibliography}{99}
\bibitem{topref}
G.\ Kasieczka, T.\ Plehn, J.\ Thompson and M.\ Russel,
\emph{Top Quark Tagging Reference Dataset},
Zenodo (2019), doi:10.5281/zenodo.2603256.

\bibitem{particlenet}
H.\ Qu and L.\ Gouskos, \emph{ParticleNet: Jet Tagging via Particle
Clouds}, Phys.\ Rev.\ D 101 (2020) 056019, arXiv:1902.08570.

\bibitem{partformer}
H.\ Qu, C.\ Li and S.\ Qian, \emph{Particle Transformer for Jet
Tagging}, ICML 2022, arXiv:2202.03772.

\bibitem{dgcnn}
Y.\ Wang \emph{et al.}, \emph{Dynamic Graph CNN for Learning on
Point Clouds}, ACM Trans.\ Graph.\ 38(5):1--12 (2019),
arXiv:1801.07829.

\bibitem{guo2017}
C.\ Guo \emph{et al.}, \emph{On Calibration of Modern Neural
Networks}, ICML 2017, arXiv:1706.04599.

\bibitem{cacciari}
M.\ Cacciari, G.\ P.\ Salam and G.\ Soyez, \emph{The anti-$k_T$ jet
clustering algorithm}, JHEP 04 (2008) 063, arXiv:0802.1189.

\bibitem{thaler}
J.\ Thaler and K.\ Van Tilburg, \emph{Identifying Boosted Objects
with N-subjettiness}, JHEP 03 (2011) 015, arXiv:1011.2268.

\bibitem{larkoski}
A.\ J.\ Larkoski, G.\ P.\ Salam and J.\ Thaler, \emph{Energy
Correlation Functions for Jet Substructure}, JHEP 06 (2013) 108,
arXiv:1305.0007.

\bibitem{pythia}
T.\ Sjöstrand \emph{et al.}, \emph{An Introduction to PYTHIA 8.2},
Comput.\ Phys.\ Commun.\ 191 (2015) 159, arXiv:1410.3012.

\bibitem{xgboost}
T.\ Chen and C.\ Guestrin, \emph{XGBoost: A Scalable Tree Boosting
System}, KDD 2016, arXiv:1603.02754.

\bibitem{kasieczka_review}
G.\ Kasieczka \emph{et al.}, \emph{The Machine Learning Landscape
of Top Taggers}, SciPost Phys.\ 7 (2019) 014, arXiv:1902.09914.
\end{thebibliography}

% ================================ APPENDICES
\appendix
\chapter{Reproducing the results}\label{app:repro}
\section{Environment}
\begin{verbatim}
python3 -m venv particle_venv
source particle_venv/bin/activate
pip install --upgrade pip
pip install -r requirements.txt
\end{verbatim}

\section{Dataset}
\begin{verbatim}
python3 download_data.py
\end{verbatim}

\section{Headline run}
\begin{verbatim}
python3 run_full_pipeline.py
\end{verbatim}

\section{Research-angle study}
\begin{verbatim}
python3 run_study.py
\end{verbatim}

\section{Compiling the documents}
\begin{verbatim}
cd docs/article   && pdflatex article.tex   && pdflatex article.tex
cd ../tutorial    && pdflatex tutorial.tex  && pdflatex tutorial.tex
cd ../thesis      && pdflatex thesis.tex    && pdflatex thesis.tex
\end{verbatim}

\chapter{Selected code excerpts}\label{app:code}
This appendix excerpts the key implementations from the project
repository.

\section{EdgeConv block (\texttt{src/particle\_net.py})}
{\footnotesize
\begin{verbatim}
def get_edge_features(x, idx):
    """
    For each point i and its neighbour j, the edge feature is
        [ x_j - x_i ,  x_i ]   (concatenation along last axis).
    """
    B, N, C = x.shape
    k = idx.shape[2]
    idx_expanded = idx.unsqueeze(-1).expand(-1,-1,-1,C)
    neighbors = torch.gather(
        x.unsqueeze(1).expand(-1,N,-1,-1), 2, idx_expanded)
    center = x.unsqueeze(2).expand(-1,-1,k,-1)
    return torch.cat([neighbors - center, center], dim=-1)
\end{verbatim}}

\section{Pairwise attention bias (\texttt{src/particle\_transformer.py})}
{\footnotesize
\begin{verbatim}
class PairwiseInteraction(nn.Module):
    def __init__(self, pair_input_dim, num_heads):
        super().__init__()
        self.mlp = nn.Sequential(
            nn.Linear(pair_input_dim, 32),
            nn.ReLU(inplace=True),
            nn.Linear(32, num_heads),
        )
    def forward(self, pair_features):
        bias = self.mlp(pair_features)         # (B,N,N,H)
        return bias.permute(0, 3, 1, 2)        # (B,H,N,N)
\end{verbatim}}

\end{document}
